\chapter{Hearing Aids}
\index{Hearing Aids}

\section*{Definition and Overview}
Hearing aids are small electronic devices worn in or behind the ear that amplify sound to help people with hearing loss hear more clearly. They work by picking up sound with a microphone, processing and amplifying it according to the wearer's specific pattern of hearing loss, and delivering it through a receiver into the ear. Modern hearing aids are digital, programmable, and highly sophisticated, with features such as directional microphones, noise reduction, and Bluetooth connectivity. Hearing aids do not restore normal hearing, but they significantly improve communication, safety, and quality of life for most people with sensorineural or mixed hearing loss.

\section*{Epidemiology}
Hearing loss is one of the most common chronic conditions worldwide, affecting roughly one in six people. Its prevalence rises steeply with age: around one in three people over 65 and one in two over 75 experience hearing loss significant enough to affect communication. Yet despite this burden, only a minority of people who could benefit from hearing aids actually use them, with uptake affected by cost, stigma, and lack of awareness. Untreated hearing loss is associated with social isolation, depression, cognitive decline, and dementia. Increasing access to hearing aids is therefore a public health priority, particularly as the population ages.

\section*{Aetiology and Pathophysiology}
\begin{itemize}
    \item \textbf{Sensorineural hearing loss:} damage to the hair cells of the cochlea or the auditory nerve, most commonly from ageing (presbycusis) or noise exposure
    \item \textbf{Conductive hearing loss:} problems conducting sound through the outer or middle ear, such as wax, middle ear infection, or otosclerosis
    \item \textbf{Mixed hearing loss:} a combination of both types
    \item \textbf{Hearing aids amplify sound to compensate:} they make speech louder and clearer for damaged ears, although they cannot repair the underlying damage
    \item \textbf{Degrees of loss:} aids are most effective for mild to moderate loss, with more severe loss often requiring additional technology such as implantable devices
\end{itemize}

\section*{Clinical Features}
\begin{itemize}
    \item Difficulty following conversations, especially in background noise
    \item Frequently asking people to repeat themselves
    \item Turning up the television or radio louder than others prefer
    \item Withdrawal from social situations because of difficulty hearing
    \item Tinnitus, ringing or buzzing in the ears, which may accompany hearing loss
    \item In children, delayed speech, language, and learning development
\end{itemize}

\section*{Differential Diagnosis}
\begin{itemize}
    \item Hearing loss should be professionally assessed to distinguish sensorineural from conductive causes
    \item A blocked ear from wax or infection may temporarily reduce hearing and is often easily treated
    \item Sudden hearing loss is a medical emergency that requires urgent treatment and differs from gradual loss
    \item Cognitive changes, including early dementia, can present with apparent difficulty understanding, and should be considered
    \item Some medications and medical conditions can cause or worsen hearing loss
\end{itemize}

\section*{When to Seek Medical Care}
Anyone who suspects hearing loss should have a hearing assessment with an audiologist or hearing service provider. Sudden hearing loss, hearing loss in one ear, hearing loss with dizziness, or hearing loss after head injury requires urgent medical review.

\textbf{Contact your local emergency services for:}
\begin{itemize}
    \item Sudden hearing loss, especially if accompanied by dizziness or ringing
    \item Hearing loss with severe headache, facial weakness, or other neurological symptoms
    \item Hearing loss following a head injury
\end{itemize}

\section*{Diagnosis and Investigations}
\begin{itemize}
    \item \textbf{Hearing test (audiometry):} measures the softest sounds heard at different frequencies to characterise the type and degree of loss
    \item \textbf{Speech testing:} assesses the ability to understand words, which guides hearing aid programming
    \item \textbf{Examination:} inspection of the ears to exclude wax, infection, or other treatable causes
    \item \textbf{Tympanometry:} tests middle ear function when conductive loss is suspected
    \item \textbf{Trial period:} most providers offer a trial so users can experience aided hearing before purchase
\end{itemize}

\section*{Management}
\begin{itemize}
    \item \textbf{Hearing aid fitting:} devices are selected and programmed to the individual's audiogram and listening needs
    \item \textbf{Styles:} behind-the-ear and in-the-ear models are available, varying in size, power, and features
    \item \textbf{Features:} directional microphones, noise reduction, feedback management, telecoils, and Bluetooth streaming improve listening in real-world situations
    \item \textbf{Auditory rehabilitation:} counselling and communication strategies help users adapt, and family members learn supportive techniques
    \item \textbf{Follow-up and adjustment:} hearing aids are fine-tuned over time as hearing changes and familiarity grows
    \item \textbf{Maintenance:} daily cleaning, battery or rechargeable care, and regular servicing extend device life
    \item \textbf{Assistance programs:} in many countries, hearing services are subsidised for eligible pensioners, veterans, and children through government programs
\end{itemize}

\section*{Complications}
\begin{itemize}
    \item Untreated hearing loss contributes to social isolation, depression, and cognitive decline
    \item Poorly fitted or adjusted aids can deliver uncomfortable or distorted sound
    \item Feedback (whistling) and occlusion effects may occur but are usually correctable
    \item Earwax impaction and skin irritation or infection can occur with poorly maintained devices
    \item Hearing aids are a financial investment, and cost can be a barrier to access
\end{itemize}

\section*{Prognosis and Outcomes}
Most people with mild to moderate hearing loss gain substantial benefit from hearing aids, with studies showing improvements in communication, social participation, mood, and quality of life. Adaptation takes time, typically weeks to months, and realistic expectations are important: aids help most in quiet situations and improve speech understanding in noise to a variable degree. Regular use is strongly associated with better outcomes, and ongoing professional support enhances satisfaction. For those with more severe loss, cochlear implants or other implantable devices offer further options.

\section*{Prevention and Self-Care}
\begin{itemize}
    \item Protect hearing from loud noise with earplugs or earmuffs, and keep personal audio volume at safe levels
    \item Have hearing checked routinely, particularly from middle age onwards
    \item Seek assessment early, because intervention is more effective when hearing loss is mild
    \item Use the trial period and involve family in the decision to trial hearing aids
    \item Wear aids consistently during the adaptation period to maximise benefit
    \item Keep devices dry, clean, and free of wax, and store them safely
    \item Manage other risk factors for hearing loss, including diabetes and cardiovascular disease
\end{itemize}

\section*{Sources and Further Reading}
The following reputable sources provide further information for healthcare students and patients seeking reliable, current advice about hearing aids.

\begin{itemize}
    \item Healthdirect. \textit{Hearing aids: how they work and who can benefit.} https://www.healthdirect.gov.au/hearing-aids
    \item Hearing Australia. \textit{Information about hearing loss and hearing aids.} https://www.hearing.com.au/
    \item Australian Government Department of Health. \textit{Hearing Services Program.} https://www.health.gov.au/
    \item NHS. \textit{Hearing aids: getting used to them.} https://www.nhs.uk/conditions/hearing-aids/
    \item Mayo Clinic. \textit{Hearing aids: how to choose and use them.} https://www.mayoclinic.org/
    \item World Health Organization. \textit{Deafness and hearing loss.} https://www.who.int/
\end{itemize}
